
En
effet, (ii) implique qu'il existe un diviseur effectif @@M0@@ tel que @@M1@@, et @@M2@@ implique que @@M3@@ ; d'o\`u (i).

Pour voir que (i') implique (ii), on peut \'evidemment supposer que @@M4@@ est tr\`es ample ; alors il existe une courbe ferm\'ee int\`egre, localement principale @@M5@@ telle que @@M6@@. D'apr\`es le th\'eor\`eme de Riemann, (i') implique que @@M7@@ pour @@M8@@ et @@M9@@ avec @@M10@@. Comme @@M11@@ pour @@M12@@, la suite exacte
@@M13@@
donne que @@M14@@ pour @@M15@@. Donc pour @@M16@@, on a
@@M17@@
ce qui ach\`eve de prouver le lemme.

\textbf{7.1.3.} Maintenant supposons dans 7.1 qu'il existe un faisceau inversible @@M18@@ sur @@M19@@ tel que @@M20@@, et raisonnons par l'absurde. Rempla\c{c}ons @@M21@@ par le sch\'ema r\'eduit sous-jacent d'une pr\'esentation de Chow de @@M22@@, ce qui est loisible gr\^ace \`a la formule de projection X 4.4.4, et soit @@M23@@ ample sur @@M24@@. Prenons @@M25@@ tel que @@M26@@ soit ample, et tel que @@M27@@ ; prenons @@M28@@ (@@M29@@) tels que, posant @@M30@@, on ait que @@M31@@ est @@M32@@ ; alors on a aussi que @@M33@@ est @@M34@@ ; ce qui revient au cas 7.1.1, avec @@M35@@ et @@M36@@.

\textbf{7.1.4.} Enfin consid\'erons le cas g\'en\'eral de 7.1. Toujours raisonnant par l'absurde, supposons que @@M37@@ ne soit pas num\'eriquement \'equivalent \`a @@M38@@ ; en outre, rempla\c{c}ons @@M39@@ par la normalis\'ee de @@M40@@, de sorte que @@M41@@ n'a 
